%---------------------------------------------------------------------
\chapter{Introduction}
%---------------------------------------------------------------------

This chapter introduces the handloom sector and the problem of verifying Geographical Indication (GI)
claims on handloom products. It states the decision that the proposed agent supports, explains why an
agentic design is needed rather than a script or a single language-model prompt, and sets out the
objectives and scope of the project.

\section{Background}

A Geographical Indication is a sign used on goods that have a specific geographical origin and that
possess qualities, a reputation or characteristics essentially attributable to that origin. The
Agreement on Trade-Related Aspects of Intellectual Property Rights (TRIPS) requires member countries
to provide legal means to prevent the use of a geographical indication on goods that do not originate
in the indicated area~\cite{trips}. In India this obligation is implemented by the Geographical
Indications of Goods (Registration and Protection) Act, 1999, under which a GI is registered by an
association of producers, and individual producers may in turn be registered as authorised
users~\cite{giact1999}. Registrations are maintained by the Geographical Indications Registry of the
Office of the Controller General of Patents, Designs and Trade Marks, commonly referred to as IP
India~\cite{ipindia}.

Handloom textiles form one of the largest groups of registered Indian GIs. Products such as
Kancheepuram Silk from Tamil Nadu, Pochampalli Ikat from Telangana, Banaras Brocades and Sarees from
Uttar Pradesh, Sambalpuri Bandha from Odisha and Chanderi Sarees from Madhya Pradesh are each tied to
named districts, weaving communities, materials and techniques~\cite{ipindia,pibgi}. The sector is
large and dispersed: according to the Fourth All India Handloom Census 2019--20, India has
26,73,891 handloom weavers and 8,48,621 allied workers, with the largest weaver populations in
Assam (11,07,428), West Bengal (3,66,656) and Manipur (2,12,481)~\cite{pibcensus}. For the three
products on which this project was first scoped, the census records 1,97,818 handloom weavers in
Tamil Nadu, 25,930 in Telangana and 1,30,778 in Uttar Pradesh~\cite{pibcensus}.

Because a GI signals both origin and traditional method, a GI claim adds commercial value to a product,
and it is used prominently in product descriptions, labels and online listings. The Ministry of
Textiles operates certification schemes such as the Handloom Mark and the India Handloom Brand to help
buyers identify genuine handloom goods~\cite{handloommark}. However, a buyer or a marketplace operator
who sees the claim ``GI-tagged Kanjivaram silk saree'' on a listing still has to establish several
separate facts before trusting it: that the product name corresponds to a registered GI, that the
stated place of manufacture lies within the registered geographical area, that the materials are
consistent with the registered description, and that the weaver and production batch can be traced to
the relevant weaving cluster.

\section{Problem Statement}

A buyer, marketplace moderator or cooperative officer who receives a handloom product with a GI-tag
claim has no single place to check whether the claim is supported. The necessary evidence is spread
across the GI registration list, the statement of case in each GI application, the national handloom
census and producer-level weaver and batch records, and the product label itself may use a local
name (``Kanjivaram'' instead of ``Kancheepuram Silk''), omit identifiers, or contain details that
contradict the registration. Checking these sources manually is slow and inconsistent, and an
unsupported claim may be accepted while a genuine product with incomplete paperwork may be wrongly
rejected. \textbf{The decision supported by this project is: is this handloom product's GI-tag claim
supported by the available official evidence, contradicted by it, or not decidable without a human
reviewer?}

\section{Motivation and Why Agentic AI}

Three simpler alternatives were considered before an agentic design was adopted.

\paragraph{A fixed rule script.} A script could compare a structured form against a lookup table.
However, product labels are free text, products are referred to by local names and spellings
(``Pochampally'', ``Ikkat'', ``Banarasi''), and the set of checks that is meaningful depends on which
fields are present: a label without a batch number cannot be checked against provenance records, and
a product that is registered for a whole state (for example, Muga Silk of Assam) must be checked
differently from one registered for specific districts. A fixed pipeline either fails on such inputs or
silently skips checks.

\paragraph{A single language-model prompt.} A large language model (LLM) asked ``is this saree
authentic?'' will answer from its training data. It cannot consult the current registry, it cannot
query producer records, and its output cannot be traced to a source. Language models are known to
produce fluent but unsupported statements~\cite{ji2023hallucination}, and smaller models are
particularly weak when relevant facts are buried in long contexts~\cite{liu2024lost}. For a decision
that affects a producer's reputation, an answer without citations is not acceptable.

\paragraph{Plain retrieval followed by one answer.} A single document lookup followed by one decision
improves grounding, but it cannot notice that its first search failed (for example, because the label
used an alias), and it has no mechanism to stop and ask a human when evidence is missing. Chapter~7
shows that such a baseline marks every case in the test set as ``Supported''.

The problem therefore requires several capabilities that together motivate an agentic design. It
requires \emph{tool use}, because facts must come from registry, census and producer records rather
than from model memory~\cite{schick2023toolformer}. It requires \emph{retrieval over official
documents}, so that every verdict can cite the registered description~\cite{lewis2020rag}. It requires
\emph{planning and re-planning}, because the checks needed differ by claim and a weak first search must
be detected and repeated with a better query~\cite{yan2024crag}. It requires \emph{evidence grading
(reflection)} before a verdict is issued, so that the agent does not guess when evidence is incomplete.
Finally, it requires \emph{human-in-the-loop (HITL) escalation}, because abstaining and referring a case
to a reviewer is safer than wrongly certifying or wrongly rejecting a product.

\section{Objectives}

The objectives of the project are as follows.

\begin{enumerate}[itemsep=2pt]
\item To design an agent, implemented as a state graph, that verifies a handloom GI claim against
official GI and census data and producer records, and returns one of three advisory verdicts with
citations.
\item To implement seven typed tools and retrieval over official GI text, with a bounded re-planning
loop (at most three iterations) that refines the search when the first attempt is weak.
\item To implement a guardrail and a human review workflow so that every case without sufficient,
cited evidence is routed to a reviewer, whose decision is recorded and visible to the requester.
\item To record a step-level trace and a persistent checkpoint for every node of every run.
\item To evaluate the agent on a labelled set of 45 cases against a single-lookup baseline, with
targets of at least 80\% verdict accuracy and at least 90\% correct abstention.
\item To measure, through a controlled A/B experiment, the effect of the re-planning loop on accuracy,
latency and tool usage.
\end{enumerate}

\section{Scope of the Project}

\paragraph{Included.} The system verifies GI-tag claims for fifteen registered handloom GIs: the three
products named in the project's design document (Kancheepuram Silk, Pochampalli Ikat, and Banaras
Brocades and Sarees) and twelve further handloom GIs whose geographical area and materials are stated
in a Government of India GI catalogue~\cite{pibgi}. Input is a product label or seller claim
containing some or all of the following fields: product name, region, state, material, weaver
identifier and batch identifier. Output is an advisory verdict with the list of checks performed and
the sources cited. The system includes a web interface for requesters, a reviewer dashboard, and an
offline evaluation harness.

\paragraph{Excluded.} The system does not issue legal certification and does not replace the GI
Registry or the authorised-user process under the GI Act~\cite{giact1999}. It does not physically
inspect textiles. It does not access live government portals at run time; reference data are stored
locally. Individual weaver and batch records are not publicly available, and the records used are
simulated, as stated in Section~3.4. Products outside the fifteen GIs are reported as ``not covered''
and may be sent to a reviewer, but they are not verified automatically. Label images are supported
only when the optional optical character recognition (OCR) package is installed, and this path was not
part of the evaluation.

\section{Organisation of the Report}

Chapter~2 reviews how GI claims are verified in practice and the technical literature on tool use,
retrieval, planning, reflection and evaluation. Chapter~3 states the stakeholders, requirements and
data sources. Chapter~4 presents the architecture of LoomTrace, its components, the agentic patterns
chosen and the control and data flow. Chapter~5 maps the system to the agentic concepts of the course
with evidence from real runs. Chapter~6 describes the implementation, including the technology stack,
the agent loop, the tools, the guardrails and the user interface. Chapter~7 describes tracing and
reports the evaluation against a baseline, and Chapter~8 reports the A/B experiment on the re-planning
loop. Chapter~9 discusses limitations, failure modes and ethical considerations, and Chapter~10
concludes with future work.
